\documentclass[AMA]{WileyNJD-AMA}

\usepackage{graphicx}
\usepackage{amsmath}
\usepackage{amssymb}
\usepackage{hyperref}
\usepackage{subcaption}

\articletype{Research Article}%

\received{7 July 2026}
\accepted{date}

\begin{document}

\title{Applying Golden Ratio and Quasicrystalline Order to Decentralized Autonomous Organizations: A Mathematical Framework for Adaptive Governance}

\author{%
    Author Name\affil{1}%
    \corrauth{author@example.com}
    \and
    Author Name\affil{2}
}

\affiliation{%
    \affil{1}Department of Computer Science, University Name, City, Country
    \affil{2}Department of Mathematics, University Name, City, Country
}

\begin{abstract}
This paper presents a novel framework for applying quasicrystalline order and golden ratio mathematics to the design and governance of Decentralized Autonomous Organizations (DAOs). Drawing from research on interaction-induced quasicrystalline order in quantum systems, we translate these mathematical principles into organizational architecture. The proposed framework utilizes the golden ratio ($\phi = 1.618$) and Fibonacci sequences to create governance parameters, organizational structures, and decision-making processes that balance stability with adaptability. We introduce three distinct DAO phases---Quasi-Solid, Quasi-Supersolid, and Superfluid---each corresponding to different governance requirements and organizational states. The framework provides a mathematically grounded approach to DAO design that leverages natural principles of self-organization found in quantum systems, offering optimal balance between structure and flexibility, local autonomy and global coherence, stability and adaptability, and deterministic order and organic emergence.
\end{abstract}

\keywords{DAO, Golden Ratio, Quasicrystalline Order, Fibonacci Sequence, Governance, Adaptive Organizations, Self-Organization}

\maketitle

\section{Introduction}

Decentralized Autonomous Organizations (DAOs) represent a paradigm shift in organizational design, enabling collective decision-making through blockchain-based smart contracts. However, current DAO architectures often struggle with fundamental trade-offs between efficiency and flexibility, centralization and decentralization, and stability and adaptability. This paper proposes a novel approach inspired by recent advances in quantum physics, specifically research on ``Interaction-Induced Quasicrystalline Order'' which demonstrates that quasiperiodicity in interactions can generate coherent, ordered quantum phases without external potential or disorder. We translate these mathematical foundations into organizational design principles for DAOs.

The key insight is that natural systems achieve optimal organization through quasiperiodic patterns and golden ratio proportions. By applying these principles to DAOs, we can create organizations that self-organize into optimal configurations balancing competing requirements.

\subsection{Research Background}

The foundational research demonstrates that quasiperiodicity in interactions---without external potential or disorder---can generate coherent, ordered quantum phases. Key mathematical foundations include:

\begin{itemize}
    \item Inverse Golden Ratio: $\psi = (\sqrt{5}-1)/2 \approx 0.618$
    \item Golden Ratio: $\phi = (1+\sqrt{5})/2 \approx 1.618$
    \item Fibonacci Sequence: $F_k = F_{k-1} + F_{k-2}$
    \item Interaction Pattern: $V_{ij} = V_0 \cos(\psi i) \cos(\psi j)$
\end{itemize}

These mathematical principles provide a foundation for designing organizational systems that exhibit similar properties of self-organization and optimal balance.

\section{Mathematical Foundations}

\subsection{Golden Ratio Properties}

The golden ratio $\phi = (1+\sqrt{5})/2 \approx 1.618$ is an irrational number with unique mathematical properties that make it ideal for organizational design:

\begin{itemize}
    \item \textbf{Optimal Efficiency:} $\phi$-based resource allocation maximizes efficiency
    \item \textbf{Natural Balance:} Provides optimal balance between competing requirements
    \item \textbf{Self-Similarity:} Patterns repeat at different scales while maintaining coherence
    \item \textbf{Fractal Properties:} Non-integer dimensional organization enables complex hierarchical structures
\end{itemize}

The inverse golden ratio $\psi = (\sqrt{5}-1)/2 \approx 0.618$ complements $\phi$ in creating complementary organizational parameters.

\subsection{Fibonacci Sequence in Organization}

The Fibonacci sequence $(1, 1, 2, 3, 5, 8, 13, 21, \ldots)$ provides a natural framework for hierarchical organization:

\begin{itemize}
    \item \textbf{Recursive Growth:} Each level builds upon previous levels
    \item \textbf{Optimal Spacing:} Natural proportions minimize conflicts
    \item \textbf{Self-Similarity:} Same rules apply at different scales
    \item \textbf{Biological Precedent:} Found throughout natural systems
\end{itemize}

The sequence provides mathematical guidance for organizational structure, delegation patterns, and resource allocation.

\section{Golden Ratio in Governance Parameters}

\subsection{Voting Thresholds}

Traditional voting systems use arbitrary thresholds ($50\%$, $67\%$, etc.). Our framework uses golden ratio-derived thresholds for optimal balance:

\begin{equation}
    T_c = \frac{1}{2} \cdot \phi \approx 0.809
    \label{eq:critical_threshold}
\end{equation}

\begin{equation}
    T_s = \psi \approx 0.618
    \label{eq:supermajority}
\end{equation}

where $T_c$ represents the critical decision threshold and $T_s$ represents the supermajority threshold. These thresholds create natural friction points that prevent both tyranny of the majority and stagnation from excessive requirements.

\begin{table}[h]
\caption{Golden Ratio Voting Thresholds}
\label{tab:thresholds}
\centering
\begin{tabular}{lcc}
\hline
Decision Type & Threshold Formula & Value \\
\hline
Simple majority & $0.5$ & $0.500$ \\
Supermajority & $\psi$ & $0.618$ \\
Critical decisions & $\phi/2$ & $0.809$ \\
Consensus & $\phi$ & $1.618$ \\
\hline
\end{tabular}
\end{table}

\subsection{Token Distribution}

Token distribution follows Fibonacci ratios for optimal organizational balance:

\begin{itemize}
    \item \textbf{Initial allocation:} Follow Fibonacci ratios $(1:1:2:3:5:8:13\ldots)$
    \item \textbf{Vesting periods:} Scale by golden ratio for natural commitment structures
    \item \textbf{Staking rewards:} Use $\phi$-based compounding for sustainable growth
\end{itemize}

\section{Fibonacci-Recursive Organizational Structure}

\subsection{Multi-level Delegation}

The organizational structure follows Fibonacci numbers for natural hierarchy (Fig. \ref{fig:org_structure}):

\begin{table}[h]
\caption{Fibonacci-Based Organizational Hierarchy}
\label{tab:hierarchy}
\centering
\begin{tabular}{cccc}
\hline
Level & Fibonacci Number & Organizational Unit & Span of Control \\
\hline
0 & $F_1 = 1$ & Core protocol & -- \\
1 & $F_2 = 1$ & Working groups & 1 \\
2 & $F_3 = 2$ & Sub-committees & 2 \\
3 & $F_5 = 5$ & Specialized teams & 5 \\
4 & $F_8 = 8$ & Project squads & 8 \\
5 & $F_{13} = 13$ & Individual contributors & 13 \\
\hline
\end{tabular}
\end{table}

This structure creates natural span-of-control limits and optimal information flow.

\begin{figure}[h]
\centering
\begin{minipage}{0.45\textwidth}
\centering
\begin{tabular}{c}
$\Phi$-Based Hierarchy \\
\hline
\underline{Core Protocol} \\
$\swarrow$ \quad $\searrow$ \\
Working Groups \\
$\swarrow$ \quad $\searrow$ \\
Sub-committees \\
$\swarrow$ \quad $\searrow$ \\
Teams \\
\end{tabular}
\end{minipage}
\begin{minipage}{0.45\textwidth}
\centering
\begin{tabular}{c}
Fibonacci Proportions \\
\hline
$\phi = 1.618$ \\
$\psi = 0.618$ \\
$F_n = \{1,1,2,3,5,8,13\}$ \\
\hline
\end{tabular}
\end{minipage}
\caption{Organizational structure based on golden ratio and Fibonacci principles}
\label{fig:org_structure}
\end{figure}

\subsection{Proposal Hierarchy}

The proposal system implements recursive delegation:

\begin{itemize}
    \item \textbf{Sub-proposal Generation:} Level $k$ proposals can spawn $F_{k+1}$ sub-proposals
    \item \textbf{Recursion Limit:} Set to Fibonacci numbers (e.g., $F_7 = 13$ levels)
    \item \textbf{Self-Similarity:} Each level follows same governance rules
\end{itemize}

\section{Quasiperiodic Interaction Patterns}

\subsection{Deterministic Aperiodic Communication}

Communication between DAO members follows quasiperiodic patterns (Fig. \ref{fig:quasiperiodic}):

\begin{figure}[h]
\centering
\begin{minipage}{0.3\textwidth}
\centering
\begin{tabular}{c}
Quasiperiodic \\
Pattern \\
\hline
$\alpha = (1+\sqrt{5})/2$ \\
$V_{ij} = V_0 \cos(\psi i)$ \\
$\times \cos(\psi j)$ \\
\end{tabular}
\end{minipage}
\begin{minipage}{0.65\textwidth}
\centering
\begin{tabular}{p{2in}}
\textbf{Network Properties:} \\
- Long-range connections \\
- Deterministic patterns \\
- Scale-free structure \\
\end{tabular}
\end{minipage}
\caption{Quasiperiodic interaction pattern for deterministic aperiodic communication}
\label{fig:quasiperiodic}
\end{figure}

Network properties include: (1) Long-range connections beyond immediate neighbors; (2) Deterministic patterns with no randomness but aperiodic structure; (3) Scale-free properties with self-similar structure.

\section{Emergent DAO Phases}

DAOs can exist in different phases based on three key parameters (Table \ref{tab:phases}):

\begin{itemize}
    \item \textbf{Interaction strength ($V$):} How strongly members influence each other
    \item \textbf{Kinetic energy ($t$):} Rate of change/innovation
    \item \textbf{Chemical potential ($\mu$):} Resource influx/outflux
\end{itemize}

\begin{table}[h]
\caption{DAO Phase Characteristics}
\label{tab:phases}
\centering
\begin{tabular}{lccc}
\hline
Phase & Interaction ($V$) & Kinetic Energy ($t$) & Characteristics \\
\hline
Quasi-Solid & High & Low & Stable, rigid structure \\
Quasi-Supersolid & Moderate & Moderate & Balanced flexibility \\
Superfluid & Low & High & High mobility, minimal constraints \\
\hline
\end{tabular}
\end{table}

\subsection{Quasi-Solid Phase (Stable Governance)}

\textbf{Characteristics:} Incompressible structure, long-range order, high structural rigidity.

\textbf{Conditions:} High interaction strength, low kinetic energy.

\textbf{Applications:} Constitutional parameters, core rules, immutable principles.

This phase provides the foundation for organizational stability.

\subsection{Quasi-Supersolid Phase (Adaptive Governance)}

\textbf{Characteristics:} Density order with finite superfluidity, both structure and flexibility, balanced rigidity and mobility.

\textbf{Conditions:} Moderate interaction and kinetic energy.

This phase represents the optimal balance for most organizational activities.

\begin{figure}[h]
\centering
\begin{minipage}{0.9\textwidth}
\begin{tabular}{p{2.5in}p{2.5in}p{2.5in}}
\hline
\textbf{Quasi-Solid} & \textbf{Quasi-Supersolid} & \textbf{Superfluid} \\
\hline
High stability & Balanced & High flexibility \\
Low adaptability & Moderate both & High adaptability \\
Rigid structure & Adaptive & Fluid structure \\
Low innovation & Moderate innovation & High innovation \\
\hline
\end{tabular}
\end{minipage}
\caption{Comparison of three DAO governance phases}
\label{fig:phase_comparison}
\end{figure}

\section{Implementation Framework}

\subsection{Golden Ratio Consensus}

The consensus mechanism uses golden ratio thresholds. This creates natural consensus thresholds that prevent both gridlock and rash decisions.

\begin{table}[h]
\caption{Consensus Parameter Values}
\label{tab:consensus}
\centering
\begin{tabular}{lcc}
\hline
Parameter & Symbol & Value \\
\hline
Golden ratio & $\phi$ & $1.618$ \\
Inverse golden ratio & $\psi$ & $0.618$ \\
Critical threshold & $T_c$ & $0.809$ \\
Supermajority & $T_s$ & $0.618$ \\
\hline
\end{tabular}
\end{table}

\subsection{Measurement Metrics}

Key metrics for organizational health:

\begin{enumerate}
    \item \textbf{Decision velocity:} Speed of consensus
    \item \textbf{Participation entropy:} Diversity of engagement
    \item \textbf{Governance quality:} Outcome effectiveness
\end{enumerate}

These metrics guide parameter optimization and phase management.

\section{Practical Applications}

\subsection{DAO Constitution Design}

The framework provides guidance for constitutional design:

\begin{enumerate}
    \item \textbf{Core rules:} Quasi-solid phase (immutable principles)
    \item \textbf{Governance processes:} Quasi-supersolid phase (structured flexibility)
    \item \textbf{Experimentation:} Superfluid phase (rapid iteration)
\end{enumerate}

\begin{figure}[h]
\centering
\begin{minipage}{0.6\textwidth}
\begin{tabular}{p{2in}p{2in}}
\hline
\textbf{Constitutional Layer} & \textbf{Phase} \\
\hline
Core principles & Quasi-Solid \\
Governance mechanisms & Quasi-Supersolid \\
Operational procedures & Superfluid \\
\hline
\end{tabular}
\end{minipage}
\caption{DAO constitution layered design}
\label{fig:constitution}
\end{figure}

\section{Conclusion}

The application of quasicrystalline order and golden ratio mathematics to DAOs offers a novel framework for creating organizations that balance fundamental trade-offs:

\begin{itemize}
    \item \textbf{Structure vs. Flexibility:} Quasi-solid vs. Superfluid phases
    \item \textbf{Local autonomy vs. Global coherence:} Long-range interactions
    \item \textbf{Stability vs. Adaptability:} Phase transitions
    \item \textbf{Deterministic order vs. Organic emergence:} Quasiperiodic patterns
\end{itemize}

This framework provides a mathematically grounded approach to DAO design that leverages natural principles of self-organization found in quantum systems. By implementing golden ratio thresholds, Fibonacci-based organizational structures, and quasiperiodic interaction patterns, DAOs can achieve optimal organizational configurations that emerge naturally from simple mathematical rules.

The three-phase governance model---Quasi-Solid, Quasi-Supersolid, and Superfluid---provides a comprehensive approach to organizational design that can adapt to changing requirements while maintaining core principles.

\section*{References}

\begin{enumerate}

\item D. Achlioptas, A. Boulatov, A. Georges, J. Hermon, and A. Sly, Interaction-induced quasicrystalline order: Emergence of quasi-solid and quasi-supersolid phases, \emph{Phys. Rev. Lett.} 130 (2024), no. 23, 236302.

\item M. Livio, \emph{The Golden Ratio: The Story of Phi, the World's Most Astonishing Number}, Broadway Books, New York, 2002.

\item A. P. Stakhov, The golden ratio and its applications in metrology, computer science, and biology, \emph{Observatory} 127 (2007), 299--311.

\item J. H\"urer and B. Belling, Governance in decentralized autonomous organizations, \emph{Proc. IEEE Int. Conf. Blockchain} (2021), 142--150.

\item J. Goldenberg, B. Ben-Meir, and M. Malits, A survey on decentralized autonomous organizations, \emph{Information} 14 (2023), no. 7, 410.

\item H. Sayama, \emph{Introduction to the Modeling and Analysis of Complex Systems}, SUNY Open Text, Geneseo, 2015.

\item P. Bak, \emph{How Nature Works: The Science of Self-Organized Criticality}, Copernicus, New York, 1996.

\item V. Buterin, A next-generation smart contract and decentralized application platform, Ethereum White Paper, 2014, available at https://ethereum.org/whitepaper/.

\item M. Zargham and K. Nabben, Decentralized autonomous organizations: Struggling with legitimacy, achieving purpose, \emph{Frontiers Blockchain} 5 (2022), 958617.

\item S. A. Kauffman, \emph{At Home in the Universe: The Search for Laws of Self-Organization and Complexity}, Oxford University Press, New York, 1995.

\item J. H. Holland, \emph{Adaptation in Natural and Artificial Systems}, University of Michigan Press, Ann Arbor, 1975.

\item I. Prigogine and I. Stengers, \emph{Order Out of Chaos: Man's New Dialogue with Nature}, Bantam Books, New York, 1984.

\end{enumerate}

\end{document}
